\documentclass[11pt]{article}

\usepackage[margin=1in]{geometry}
\usepackage[T1]{fontenc}
\usepackage{lmodern}
\usepackage{microtype}
\usepackage{amsmath}
\usepackage{booktabs}
\usepackage{tabularx}
\usepackage{enumitem}
\usepackage[dvipsnames]{xcolor}
\usepackage{caption}
\usepackage{xurl}
\definecolor{linkblue}{HTML}{1A5FB4}
\usepackage[colorlinks=true,citecolor=linkblue,linkcolor=linkblue,urlcolor=linkblue]{hyperref}

\setlength{\parskip}{0.45em}
\setlength{\parindent}{0pt}
\setlist{itemsep=0.25em,topsep=0.3em}
\captionsetup{font=small,labelfont=bf}
\newcolumntype{L}{>{\raggedright\arraybackslash}X}
\newcolumntype{R}{>{\raggedleft\arraybackslash}X}

\title{\textbf{Critical Minerals State Monitor: Methods Note}\\[0.3em]\large Release 1.0}
\author{Samson M'boueke}
\date{November 18, 2026}

\begin{document}
\maketitle

\begin{abstract}
\noindent The Critical Minerals State Monitor measures how exposed each of the fifty U.S.\
states and the District of Columbia is to world price movements in critical minerals. It
measures exposure on two sides, kept separate: as a producer, through the value of the
state's mine output, and as a user, through the metal its industries buy. Both are reported
in dollars and as a share of state GDP, together with the first-round net effect of a 10
percent rise in world prices. Release 1.0 covers six metals on the final 2025 U.S.\ List of
Critical Minerals that have free monthly world prices: aluminum, copper, lead, nickel, tin
and zinc. This note documents the definitions, data and limitations of release 1.0. Every
figure is rebuilt from public data by the code in the Monitor's repository.
\end{abstract}

\section{Measures}

For state $s$ and metal $m$, the Monitor reports three quantities.

\paragraph{Producer exposure} is the value of the state's mine output at the current world
price,
\begin{equation}
  P_{s,m} \;=\; Q_{s,m} \times \bar p_m ,
  \label{eq:producer}
\end{equation}
where $Q_{s,m}$ is the state's mine output in metric tons of metal content, as published by
the U.S.\ Geological Survey (USGS), and $\bar p_m$ is the average world price over the latest
twelve months.

\paragraph{User exposure} is the value of the metal bought by the state's industries outside
the mining and metal-processing chain, at the current world price:
\begin{equation}
  U_{s,m} \;=\; C_m \,\bar p_m \times
  \frac{w_{g(m),s}}{\sum_{s'} w_{g(m),s'}},
  \qquad
  w_{g,s} \;=\; \sum_{L} a_{g,L}\, X_{L,s} .
  \label{eq:user}
\end{equation}
The national level is $C_m \bar p_m$, U.S.\ apparent consumption of the metal valued at the
current price. It is shared across states in proportion to $w_{g,s}$, the state's purchases of
the input-output commodity group $g(m)$ that contains the metal. Here $a_{g,L}$ is the
national purchase of group $g$ per dollar of output in industry $L$, and $X_{L,s}$ is the
state's real output in industry $L$. Section~\ref{sec:user} defines each term.

\paragraph{Net effect} of a 10 percent rise in the world price, as a percent of state GDP, is
\begin{equation}
  N_{s,m} \;=\; 0.1 \times \frac{P_{s,m} - U_{s,m}}{Y_s} \times 100 ,
  \label{eq:net}
\end{equation}
where $Y_s$ is current-dollar state GDP. The all-metals figure sums $N_{s,m}$ over the six
metals, which assumes all six prices rise by 10 percent together.

Keeping the producer and user sides separate follows the author's job market paper, which
measures commodity price shocks through separate export and import price indices with
predetermined weights. A state can be heavily exposed on both sides at once, and the two
sides matter to different people: mine owners, workers and severance tax revenues on one
side, metal-using manufacturers and their workers on the other.

\section{Prices}

Prices are the World Bank's monthly Commodity Price Data (the ``Pink Sheet''), in US dollars
per metric ton. The current price $\bar p_m$ is the average of the latest twelve months; in
release 1.0 that is October 2025 to September 2026. Averaging over a year matches the annual
quantities it multiplies and removes month-to-month noise. Table~\ref{tab:prices} reports the
prices used and their change from the previous twelve-month average.

\begin{table}[ht]
\centering
\small
\caption{World prices used in release 1.0}
\label{tab:prices}
\begin{tabularx}{0.82\textwidth}{@{}L R R@{}}
\toprule
Metal & Average price, Oct 2025 to Sep 2026 (dollars per metric ton)
  & Change from previous 12 months (percent) \\
\midrule
Aluminum & 3,206  & $+24.9$ \\
Copper   & 12,852 & $+35.8$ \\
Lead     & 1,928  & $-2.2$ \\
Nickel   & 16,729 & $+8.4$ \\
Tin      & 48,102 & $+49.9$ \\
Zinc     & 3,425  & $+20.7$ \\
\bottomrule
\end{tabularx}
\par\smallskip
\begin{minipage}{0.82\textwidth}\footnotesize
Source: World Bank, Commodity Price Data (Pink Sheet), monthly; author's calculations.
\end{minipage}
\end{table}

\section{Producer exposure}

USGS publishes U.S.\ mine production by state for few metals, and it withholds a state's
figure where publishing it would disclose an individual company's output. The Monitor uses
the most detailed figure USGS publishes for each metal and never imputes a withheld one.
Table~\ref{tab:producer} lists the source and coverage for each metal.

\begin{table}[ht]
\centering
\small
\caption{Producer data by metal, release 1.0}
\label{tab:producer}
\begin{tabularx}{\textwidth}{@{}>{\raggedright\arraybackslash}p{1.7cm} L L@{}}
\toprule
Metal & State figures used & Withheld or not covered \\
\midrule
Copper & Output of the leading mines in 2022 (Minerals Yearbook 2022, copper, table 2),
  summed by the state of each mine: Arizona 857, Utah 179, New Mexico 86, Nevada 61 and
  Michigan 16 thousand metric tons. These mines account for 1,199 of the 1,230 thousand
  tons mined in the United States in 2022.
  & Montana and Missouri, whose output USGS withholds; the remainder of national output
  not listed by mine. \\
\addlinespace
Zinc & Alaska, 598 thousand metric tons of recoverable zinc in 2022 (Minerals Yearbook
  2022, zinc, table 2), out of 747 thousand nationally.
  & Idaho, Missouri, New York and Tennessee, reported by USGS only as a combined total. \\
\addlinespace
Nickel & Michigan, 7,490 metric tons in 2024 (Mineral Commodity Summaries 2026). The Eagle
  mine in Michigan is the only U.S.\ nickel mine, so the national mine figure is
  Michigan's.
  & None. \\
\addlinespace
Lead & None.
  & Missouri, Alaska and Idaho produce lead; USGS withholds the state values. \\
\addlinespace
Tin & None. & The United States has no tin mine production. \\
\addlinespace
Aluminum & None.
  & Primary aluminum output by state is not compiled in this release. \\
\bottomrule
\end{tabularx}
\end{table}

Quantities are from the latest year USGS publishes by state or mine (2022 for copper and
zinc, 2024 for nickel) and are valued at the current price. Withheld states are flagged in
the data files and marked on the Monitor's map; their producer exposure is reported as zero,
so it understates their true exposure.

\section{User exposure}\label{sec:user}

\paragraph{National level.} The national total for each metal, $C_m$, is U.S.\ apparent
consumption in 2024 from the USGS Mineral Commodity Summaries 2026, in metric tons of metal
content: aluminum 5.83 million, copper 1.86 million (primary refined copper plus old scrap),
lead 1.41 million, nickel 211,300 (apparent primary consumption of 90,300 plus 121,000 of
purchased scrap), tin 34,600 and zinc 808,000. Valued at current prices, these give national
user totals of about 18.7 billion dollars for aluminum, 23.9 billion for copper, 2.7 billion
for lead, 3.5 billion for nickel, 1.7 billion for tin and 2.8 billion for zinc.

The level is taken from USGS rather than from input-output purchase values because the
input-output commodities are broader than the metal itself: their values include fabrication
margins and, for nonferrous refining, gold and silver. Using purchase values as the level
overstated metal use several times over, by about five times for nickel.

\paragraph{Allocation across states.} The national total is shared across states with the
Bureau of Economic Analysis (BEA) 2017 detail Use table (before redefinitions, at producers'
prices). Each metal is matched to the input-output commodities that carry it
(Table~\ref{tab:groups}). For each industry, $a_{g,L}$ is the 2017 purchase of those
commodities per dollar of 2017 industry output. Purchases by the mining and metal-processing
industries themselves (the five commodities in Table~\ref{tab:groups}) are excluded, because
those industries pass the metal on rather than use it; the user side counts the industries
that turn metal into products. Detail industries are mapped to the industry lines of BEA's
GDP by state by their NAICS codes. Purchases by industries outside those lines, mainly
government, are dropped: 0.7 percent of the aluminum group's purchases, 1.4 percent of
copper's and 1.4 percent of the nickel, lead, zinc and tin group's.

\begin{table}[ht]
\centering
\small
\caption{Input-output commodities used to allocate each metal}
\label{tab:groups}
\begin{tabularx}{\textwidth}{@{}>{\raggedright\arraybackslash}p{3.0cm} >{\raggedright\arraybackslash}p{1.6cm} L@{}}
\toprule
Metal & BEA code & Commodity \\
\midrule
Aluminum & 331313 & Alumina refining and primary aluminum production \\
         & 33131B & Aluminum product manufacturing from purchased aluminum \\
\addlinespace
Copper   & 331420 & Copper rolling, drawing, extruding and alloying \\
\addlinespace
Nickel, lead, zinc, tin & 212230 & Copper, nickel, lead, and zinc mining \\
         & 331410 & Nonferrous metal (except aluminum) smelting and refining \\
\bottomrule
\end{tabularx}
\par\smallskip
\begin{minipage}{\textwidth}\footnotesize
Source: BEA, 2017 detail Use table, before redefinitions, producers' prices. Refined copper
(within 331410) is bought mostly by wire-rod and brass mills, which belong to 331420; the
Monitor therefore allocates copper by purchases of 331420's products.
\end{minipage}
\end{table}

BEA does not split its nonferrous smelting and mining commodities by metal, so nickel, lead,
zinc and tin share one allocation across states in this release: each state's share of the
national total is the same for all four, and only the national totals differ.

\paragraph{State output.} $X_{L,s}$ is the state's real value added in industry line $L$ in
2025, in chained 2017 dollars (BEA table SAGDP9), multiplied by the national ratio of output
to value added in that industry in 2017. This converts value added to output at 2017 prices,
on the same basis as the 2017 purchase coefficients. BEA suppresses five state-industry
values for confidentiality; they are treated as zero.

\section{State GDP}

The denominator $Y_s$ is current-dollar GDP by state for 2025 (BEA table SAGDP2), the latest
year published for all states when release 1.0 was built.

\section{Results in release 1.0}

The figures move with prices and are given here only to show what the measure produces. A
uniform 10 percent rise in all six world prices raises income, in this first-round
accounting, in five states: Alaska (by 0.28 percent of state GDP), Arizona (0.17), New
Mexico (0.07), Utah (0.06) and Nevada (0.02). It lowers it most in Indiana ($-0.06$),
Alabama, Michigan and Kentucky (each about $-0.05$) and Wisconsin ($-0.03$). In these
states the largest users of the six metals are motor vehicle and parts, primary metal,
fabricated metal and electrical equipment manufacturing. In Michigan, motor vehicle and parts manufacturing accounts for about 2.4
billion of the roughly 3.8 billion dollars of the six metals used in the state.

\section{What the measures are and are not}

The net effect is an accounting measure of the first-round change in income from a price
change, holding quantities fixed. It is not an estimate of final outcomes. In particular:
\begin{itemize}
  \item \textbf{Ownership.} Producer exposure is the value of output located in the state.
    Mine profits accrue to shareholders who may live anywhere; the state keeps wages,
    local purchases and severance and income taxes. The share of a price gain that stays in
    the state is not measured in this release.
  \item \textbf{Pass-through.} User exposure assumes the world price change reaches the
    users of the metal in full. Contract pricing, hedging, inventories and the ability to
    pass costs on to customers change how much is borne and when. Estimating that
    pass-through is the subject of the Monitor's first working paper, and its results will
    be added to the Monitor when they exist.
  \item \textbf{Direct purchases only.} User exposure counts metal bought directly by each
    industry. Metal embodied in parts bought from other industries, other states or abroad
    is not traced, and the 2017 coefficients assume every state's industry uses metal as
    the national industry did in 2017.
  \item \textbf{Withheld data.} Producer exposure is understated wherever USGS withholds a
    state's figure (Table~\ref{tab:producer}), and primary aluminum producers are not yet
    covered.
  \item \textbf{Shared allocation.} Nickel, lead, zinc and tin share one state allocation.
  \item \textbf{Vintages.} Mine quantities are from 2022 or 2024, consumption from 2024,
    input-output coefficients from 2017 and state output from 2025; all are valued at
    current prices.
  \item \textbf{Validation.} The measures have not yet been validated against state
    outcomes.
\end{itemize}

\section{Data files and reproduction}

The Monitor publishes three files:
\begin{itemize}
  \item \texttt{minerals\_state.csv}: producer and user exposure in dollars and percent of
    GDP, the producer data status and the net effect, by state and metal;
  \item \texttt{minerals\_state\_industry.csv}: user exposure by state, metal and industry
    line;
  \item \texttt{minerals\_meta.json}: prices, data years and build diagnostics.
\end{itemize}
The code is at \url{https://github.com/sammbk/critical-minerals-state-monitor}, release v1.0.0:
\texttt{code/01\_fetch.py} downloads the public inputs, \texttt{code/05\_minerals.py} builds
the measures and \texttt{code/06\_build\_monitor.py} builds the site, published at
\url{https://sammbk.github.io/critical-minerals-state-monitor/}. A few USGS files must be
saved by hand because the USGS website does not serve them to scripts; the repository's
README lists them.

\section{Planned changes}

Release 1.1 adds potash and phosphate rock, which are on the critical minerals list, are
priced monthly by the World Bank, and reach users mainly through fertilizer; it also adds
primary aluminum output by state. Later releases add estimates of price pass-through, a
calculator for the state incidence of price floors and stockpile purchases and releases,
and, if a license is funded, lithium, cobalt, graphite and rare earths, whose prices are
largely sold under license.

\section*{Sources}

\begin{itemize}[leftmargin=1.2em]
  \item U.S.\ Geological Survey. Minerals Yearbook 2022: Copper; Zinc. Tables 1 and 2.
  \item U.S.\ Geological Survey. Mineral Commodity Summaries 2026: Aluminum, Copper, Lead,
    Nickel, Tin, Zinc. Salient statistics.
  \item U.S.\ Bureau of Economic Analysis. Input-Output Accounts, 2017 detail Use table,
    before redefinitions, producers' prices.
  \item U.S.\ Bureau of Economic Analysis. GDP by State: SAGDP2 (GDP in current dollars) and
    SAGDP9 (real GDP in chained 2017 dollars), by industry.
  \item World Bank. Commodity Price Data (Pink Sheet), monthly.
  \item U.S.\ Geological Survey. Final 2025 List of Critical Minerals. Federal Register,
    November 7, 2025.
\end{itemize}

\end{document}
